\chapter{Robustness Validity: Dose-Response, Targeted Occlusion, and a Size-Controlled Centerpiece (Phase 2.4)}
\label{chap:p2-4}

\section{Goal}

The pilot tested robustness at a single magnitude per perturbation, occluded with an always-centered
square, and measured explanation drift with IoU. The plan (item 4.4) asks for four corrections: sweep
each perturbation across magnitudes to get a dose-response curve; separate \emph{targeted} occlusion
(cover the object box) from \emph{random-location} occlusion, since a centered square conflates
``covered the object'' with ``disrupted the scene''; report drift and box-disappearance jointly rather
than letting vanished boxes drop out of the IoU denominator; and --- the centerpiece --- re-examine the
pilot's headline finding that ``28.3\% of stable answers hide a drifted explanation'' under a size
control, because IoU is size-biased.

\section{What Was Done}

The perturbations were already magnitude-parametrized; \texttt{occlude} gained a \texttt{location}
argument (centered, or a seeded random position), and targeted occlusion reuses
\texttt{mask\_region} on the object's own box. The robustness metric gained a size-invariant
\texttt{object\_centroid\_drift} (normalized centroid distance, shared with the stability metric) and an
explicit \texttt{object\_disappeared} flag, so a vanished box registers as an event rather than a NaN.
Script~09 gained a \texttt{-{}-severity-sweep} mode running the full grid --- blur \{4, 8, 16\}, gamma
\{1.5, 2.5, 4.0\}, contrast \{0.15, 0.3, 0.5\}, occlusion \{0.1, 0.2, 0.35\} at both centered and
random locations, plus targeted occlusion --- and writing \texttt{robustness\_sweep.csv} (2{,}603
rows), leaving the frozen \texttt{robustness.csv} intact. Nine unit tests cover the new occlusion
locations and the metric's drift/disappearance fields.

\section{Results}

\subsection{Dose-response is monotonic and legible}

\begin{center}
\begin{tabular}{p{3.2cm}p{2.0cm}p{2.4cm}p{2.4cm}p{2.2cm}}
\toprule
\textbf{perturbation} & \textbf{severity} & \textbf{answer flip} & \textbf{centroid drift} & \textbf{vanish} \\
\midrule
gamma            & 1.5 / 2.5 / 4.0    & 5.5 / 11.7 / 20.2\% & 0.011 / 0.030 / 0.090 & 0\% \\
blur             & 4 / 8 / 16         & 27.0 / 27.0 / 28.8\% & 0.118 / 0.165 / 0.175 & 0\% \\
contrast         & 0.15 / 0.3 / 0.5   & 16.6 / 10.4 / 6.1\%  & 0.048 / 0.030 / 0.016 & 0\% \\
occlude (center) & 0.1 / 0.2 / 0.35   & 15.3 / 28.2 / 33.1\% & 0.059 / 0.105 / 0.173 & 2.5--5.5\% \\
occlude (random) & 0.1 / 0.2 / 0.35   & 12.3 / 14.7 / 24.5\% & 0.048 / 0.073 / 0.143 & 3.7--6.7\% \\
\bottomrule
\end{tabular}
\captionof{table}{Dose-response over the severity grids. Answer-flip rate and size-invariant centroid
drift rise monotonically with severity for every perturbation (contrast's factor runs the other way ---
0.15 is the strongest flattening --- so its curve correctly decreases). The single fixed magnitude the
pilot reported was one point on each of these curves.}
\label{tab:p24-dose}
\end{center}

\subsection{Where you occlude matters more than how much}

Separating occlusion by location produces the phase's cleanest result:

\begin{center}
\begin{tabular}{p{6.2cm}p{3.4cm}}
\toprule
\textbf{occlusion type} & \textbf{answer-flip rate} \\
\midrule
targeted (covers the object's own box)   & \textbf{39.2\%} \\
centered square (pilot's method)         & 25.6\% \\
random location                          & \textbf{17.2\%} \\
\bottomrule
\end{tabular}
\captionof{table}{Occlusion answer-flip rate by location. Covering the object itself flips the answer
more than twice as often as a random-location occlusion of the same area; the pilot's centered square
lands in between, confirming it conflated the two.}
\label{tab:p24-occ}
\end{center}

Targeted occlusion --- actually covering the safety object --- flips the answer 39.2\% of the time,
against 17.2\% for a random-location square of matched area. The pilot's centered square (25.6\%) sits
between them because construction objects often sit near the frame center, so it partly-covered the
object some of the time. Reporting a single centered-occlusion number therefore mixed a genuine
``removing the object changes the answer'' signal with an ``occluding anything disrupts grounding''
signal. The split makes the meaningful measurement --- targeted --- explicit.

\subsection{The 28.3\% centerpiece is roughly half IoU size artifact}

The pilot's headline was that 28.3\% of stable-answer samples have \texttt{object\_box\_iou} $< 0.3$ ---
the answer held but the explanation box ``drifted.'' Reproduced here it is 23--25\%. Two controls show
much of it is IoU's size bias, not real relocation. First, bucketing stable answers by object-box area:

\begin{center}
\begin{tabular}{p{3.6cm}p{3.6cm}p{3.0cm}}
\toprule
\textbf{object-box size} & \textbf{\% with IoU $<0.3$} & \textbf{mean IoU} \\
\midrule
small  & 31.0\% & 0.582 \\
medium & 24.4\% & 0.696 \\
large  & 13.5\% & 0.800 \\
\bottomrule
\end{tabular}
\captionof{table}{The ``drift'' finding by object-box size. Small boxes register IoU-drift more than
twice as often as large ones --- the signature of IoU's size bias, not a property of the model.}
\label{tab:p24-size}
\end{center}

Second, the size-invariant centroid drift on the same stable-answer rows: IoU flags 25.0\% as drifted
($<0.3$), but the mean normalized centroid drift is only 0.062, and just 13.7\% of these boxes actually
moved by more than 0.2 of the image diagonal. So of the quarter of stable answers IoU calls
``drifted,'' only about half correspond to a real positional move; the rest are small boxes jittering
in extent around the same location, which IoU --- size-biased --- converts into a large apparent overlap
loss. The finding survives, but at roughly half its headline magnitude once size is controlled, and it
should be stated as such.

\subsection{Vanished boxes are now counted}

Occlusion makes the object box disappear on 2.5--6.7\% of reruns (targeted: 5.1\%). Under the pilot's
IoU-only reporting these produced NaN and left the drift average; they are now recorded as an explicit
disappearance rate alongside the drift, so a perturbation that erases the explanation entirely counts
as the maximal instability it is rather than vanishing from the denominator.

\section{Standard vs. Non-Standard Choices}

\textbf{Standard.} Dose-response curves, controlling a confound by stratification, and separating a
targeted from a random intervention are ordinary experimental practice.

\textbf{Non-standard, and the contribution.} The phase re-examines a headline finding of the pilot
under two independent size controls and reports that it is roughly half artifact --- rather than
carrying the $28.3\%$ forward unqualified. And it resolves the occlusion measurement into its two
constituent signals (object removal vs.\ scene disruption), which turns a single ambiguous number into
two interpretable ones. Both are cases of trusting a measurement less until it has been decomposed.

\section{Implications}

Robustness now reports dose-response curves, a targeted-vs-random occlusion split whose targeted arm is
the faithful ``does removing the object change the answer'' measure, a size-invariant drift that
halves the overstated centerpiece, and an explicit disappearance rate. The size-control lesson now
recurs across three metrics (2.2 sparsity, 2.3 stability, 2.4 robustness), which is the empirical case
for making size-invariant measures and size-band reporting standard for every box-based quantity in the
scale run. Attention-heatmap drift as a second explanation-level signal is deferred to Phase~4, where
the attribution stream it requires is built.
